\subsection{Upgrading OpenHPC Packages}  \label{appendix:upgrade}

As newer \OHPC{} releases are made available, users are encouraged to upgrade
their locally installed packages against the latest repository versions to
obtain access to bug fixes and newer component versions. This can be
accomplished with the underlying package manager as \OHPC{} packaging maintains
versioning state across releases.

\begin{enumerate*}
\item (Optional) Ensure repo metadata is current on the SMS node and in the
compute image.

\begin{lstlisting}[language=bash,keywords={}]
[sms](*\#*)  (*\clean*)
[sms](*\#*)  (*\chrootclean*)
\end{lstlisting}

\item Upgrade master (SMS) node

\begin{lstlisting}[language=bash,keywords={}]
[sms](*\#*)  (*\upgrade*) "*-ohpc"
[sms](*\#*)  (*\upgrade*) "ohpc-base"
\end{lstlisting}

\item Upgrade packages in compute image

\begin{lstlisting}[language=bash,keywords={}]
[sms](*\#*)  (*\chrootupgrade*) "*-ohpc"
[sms](*\#*)  (*\chrootupgrade*) "ohpc-base-compute"
\end{lstlisting}

\item Rebuild image(s)

\begin{lstlisting}[language=bash,keywords={},literate={BOSVER}{\baseos{}}1]
[sms](*\#*) wwctl image build BOSVER
\end{lstlisting}
\end{enumerate*}

\noindent In the case where packages were upgraded within the compute image,
you will need to reboot the compute nodes when convenient to enable the changes.
